\documentclass[aspectratio=169,10pt]{beamer}

% ------------------------------------------------------------
% Plantilla de curso: Aprendizaje por Refuerzo
% Compilar con pdfLaTeX. No requiere archivos externos.
% ------------------------------------------------------------

\usepackage[T1]{fontenc}
\usepackage[utf8]{inputenc}
\usepackage[spanish,es-nodecimaldot]{babel}
\usepackage{lmodern}
\usepackage{microtype}
\usepackage{amsmath,amssymb,mathtools}
\usepackage{booktabs}
\usepackage{listings}
\usepackage{tikz}
\usetikzlibrary{arrows.meta,calc,positioning,shapes.geometric}

% ---------- Paleta de alto contraste para proyector y pantalla ----------
% El lienzo es siempre blanco. Los tonos claros se reservan para superficies
% pequeñas; texto, líneas y títulos usan colores suficientemente oscuros.
\definecolor{Ink}{HTML}{102A43}       % azul tinta: texto principal
\definecolor{Paper}{HTML}{FFFFFF}     % fondo obligatorio de todas las diapositivas
\definecolor{Mist}{HTML}{D9E2EC}      % gris azulado para guías y divisores
\definecolor{Teal}{HTML}{005EA8}      % azul accesible: color primario
\definecolor{TealDark}{HTML}{003E73}  % azul profundo
\definecolor{Coral}{HTML}{D9480F}     % naranja oscuro: énfasis
\definecolor{Gold}{HTML}{B77900}      % ámbar oscuro: recompensa
\definecolor{Leaf}{HTML}{2F6F44}      % verde oscuro: resultados positivos
\definecolor{CodeBg}{HTML}{102A43}    % código con contraste alto
\definecolor{SoftText}{HTML}{486581}  % texto secundario legible

% ---------- Identidad visual ----------
\usefonttheme{professionalfonts}
\setbeamercolor{normal text}{fg=Ink,bg=Paper}
\setbeamercolor{structure}{fg=Teal}
\setbeamercolor{frametitle}{fg=Ink,bg=Paper}
\setbeamercolor{alerted text}{fg=Coral}
\setbeamercolor{itemize item}{fg=Coral}
\setbeamercolor{itemize subitem}{fg=Teal}
\setbeamercolor{block title}{fg=white,bg=Teal}
\setbeamercolor{block body}{fg=Ink,bg=Teal!8}
\setbeamercolor{block title alerted}{fg=white,bg=Coral}
\setbeamercolor{block body alerted}{fg=Ink,bg=Coral!8}
\setbeamercolor{block title example}{fg=white,bg=Leaf}
\setbeamercolor{block body example}{fg=Ink,bg=Leaf!9}
\setbeamercolor{background canvas}{bg=Paper}

\setbeamertemplate{navigation symbols}{}
\setbeamertemplate{blocks}[rounded][shadow=false]
\setbeamertemplate{itemize item}{\raise0.15ex\hbox{\tikz\fill[Coral] (0,0) circle (0.55ex);}}
\setbeamertemplate{itemize subitem}{\raise0.15ex\hbox{\tikz\fill[Teal] (0,0) circle (0.4ex);}}
\setlength{\leftmargini}{1.5em}

% Barra superior del título de cada diapositiva
\setbeamertemplate{frametitle}{%
  \nointerlineskip
  \begin{beamercolorbox}[wd=\paperwidth,ht=0.95cm,dp=0.25cm,leftskip=0.55cm,rightskip=0.4cm]{frametitle}
    \usebeamerfont{frametitle}\insertframetitle
    \hfill
    \tikz[baseline=-0.5ex]{
      \fill[Teal] (0,0) circle (0.10);
      \fill[Coral] (0.28,0) circle (0.10);
      \fill[Gold] (0.56,0) circle (0.10);
    }
  \end{beamercolorbox}
  \vspace{-0.12cm}
  \hspace{0.55cm}\color{Teal}\rule{1.25cm}{1.4pt}%
}

% Pie con sección y progreso visual
\newcommand{\progressbar}{%
  \begin{tikzpicture}[baseline]
    \fill[Mist] (0,0) rectangle (3.0,0.055);
    \pgfmathsetmacro{\progress}{3.0*\insertframenumber/max(1,\inserttotalframenumber)}
    \fill[Teal] (0,0) rectangle (\progress,0.055);
  \end{tikzpicture}%
}
\setbeamertemplate{footline}{%
  \leavevmode
  \begin{beamercolorbox}[wd=\paperwidth,ht=0.42cm,dp=0.14cm,leftskip=0.55cm,rightskip=0.55cm]{normal text}
    \scriptsize\color{SoftText}\insertsectionhead
    \hfill\progressbar\hspace{0.6em}\insertframenumber
  \end{beamercolorbox}
}

% ---------- Comandos propios del curso ----------
\newcommand{\state}[1]{\tikz[baseline=(x.base)]\node[rounded corners=2pt,fill=Teal!13,text=TealDark,inner xsep=5pt,inner ysep=2pt] (x) {\scriptsize\bfseries ESTADO\enspace #1};}
\newcommand{\rlaction}[1]{\tikz[baseline=(x.base)]\node[rounded corners=2pt,fill=Coral!13,text=Coral,inner xsep=5pt,inner ysep=2pt] (x) {\scriptsize\bfseries ACCIÓN\enspace #1};}
\newcommand{\reward}[1]{\tikz[baseline=(x.base)]\node[rounded corners=2pt,fill=Gold!22,text=Ink,inner xsep=5pt,inner ysep=2pt] (x) {\scriptsize\bfseries RECOMPENSA\enspace #1};}
\newcommand{\keyidea}[1]{%
  \begin{beamercolorbox}[rounded=true,sep=0.65em,wd=\linewidth]{block body alerted}
    \textcolor{Coral}{\bfseries Idea clave}\quad #1
  \end{beamercolorbox}%
}
\newcommand{\module}[2]{%
  \begin{frame}[plain]
    \begin{tikzpicture}[remember picture,overlay]
      \fill[Paper] (current page.south west) rectangle (current page.north east);
      \fill[Teal] (current page.south west) rectangle ([xshift=0.085\paperwidth]current page.north west);
      \fill[Coral] ([xshift=0.085\paperwidth,yshift=0.8cm]current page.south west)
        rectangle ([xshift=0.097\paperwidth,yshift=-0.8cm]current page.north west);
      \draw[Mist,line width=1pt] ([xshift=0.15\paperwidth]current page.west)
        -- ([xshift=-0.8cm]current page.east);
      \node[anchor=west,text=Ink,font=\fontsize{26}{29}\selectfont\bfseries,
            text width=0.72\paperwidth] at ([xshift=0.17\paperwidth,yshift=0.45cm]current page.west) {#1};
      \node[anchor=west,text=SoftText,font=\large,text width=0.68\paperwidth]
        at ([xshift=0.17\paperwidth,yshift=-1.0cm]current page.west) {#2};
      \node[rotate=90,text=white,font=\bfseries\small] at
        ([xshift=0.0425\paperwidth]current page.west) {MÓDULO};
      \fill[Teal] ([xshift=-1.55cm,yshift=0.38cm]current page.east) circle (0.10cm);
      \fill[Coral] ([xshift=-1.22cm,yshift=0.38cm]current page.east) circle (0.10cm);
      \fill[Gold] ([xshift=-0.89cm,yshift=0.38cm]current page.east) circle (0.10cm);
    \end{tikzpicture}
  \end{frame}%
}

% ---------- Código ----------
\lstdefinestyle{rlcode}{
  language=Python,
  backgroundcolor=\color{CodeBg},
  basicstyle=\ttfamily\scriptsize\color{white!88},
  keywordstyle=\bfseries\color{Gold},
  commentstyle=\itshape\color{Leaf!85},
  stringstyle=\color{Coral!85},
  numbers=left,
  numberstyle=\tiny\color{white!35},
  numbersep=8pt,
  frame=none,
  showstringspaces=false,
  breaklines=true,
  tabsize=2,
  xleftmargin=1.3em,
  aboveskip=0pt,
  belowskip=0pt
}

% ---------- Metadatos: personaliza aquí ----------
\title{Aprendizaje por Refuerzo}
\subtitle{De la interacción a las decisiones inteligentes}
\author{Nombre del docente}
\institute{Universidad o institución · Curso 2026}
\date{}

\begin{document}

% ============================================================
% PORTADA
% ============================================================
\begin{frame}[plain]
  \begin{tikzpicture}[remember picture,overlay]
    \fill[Paper] (current page.south west) rectangle (current page.north east);
    \draw[Ink,line width=1.4pt,rounded corners=7pt]
      ([xshift=0.72\paperwidth,yshift=0.7cm]current page.south west)
      rectangle ([xshift=-0.7cm,yshift=-0.7cm]current page.north east);
    \fill[Teal] ([xshift=0.68\paperwidth]current page.south west)
      -- ([xshift=0.74\paperwidth]current page.south west)
      -- ([xshift=0.68\paperwidth,yshift=0.25\paperheight]current page.south west) -- cycle;

    % Trayectoria abstracta de un agente
    \draw[SoftText!55,line width=1.2pt,dashed]
      ([xshift=-3.6cm,yshift=-2.3cm]current page.north east)
      .. controls +(0.2,-1.1) and +(-1.0,0.5) ..
      ([xshift=-1.2cm,yshift=-4.2cm]current page.north east)
      .. controls +(0.4,-0.8) and +(0.5,0.8) ..
      ([xshift=-2.4cm,yshift=-6.4cm]current page.north east);
    \foreach \x/\y/\c in {-3.6/-2.3/Teal,-1.2/-4.2/Coral,-2.4/-6.4/Gold}{
      \fill[\c] ([xshift=\x cm,yshift=\y cm]current page.north east) circle (0.13cm);
      \draw[Ink,line width=1pt] ([xshift=\x cm,yshift=\y cm]current page.north east) circle (0.23cm);
    }

    \node[anchor=west,text=Teal,font=\bfseries\small] at ([xshift=0.75cm,yshift=2.55cm]current page.west)
      {CURSO · TEORÍA + PRÁCTICA};
    \node[anchor=west,text=Ink,font=\fontsize{28}{31}\selectfont\bfseries,text width=0.58\paperwidth]
      at ([xshift=0.75cm,yshift=1.15cm]current page.west) {\inserttitle};
    \node[anchor=west,text=SoftText,font=\large,text width=0.55\paperwidth]
      at ([xshift=0.75cm,yshift=-0.15cm]current page.west) {\insertsubtitle};
    \node[anchor=west,text=Ink,font=\small]
      at ([xshift=0.75cm,yshift=-2.35cm]current page.west) {\textbf{\insertauthor}\quad\textcolor{SoftText}{\insertinstitute}};
  \end{tikzpicture}
\end{frame}

% ============================================================
\section{Mapa del curso}
% ============================================================

\begin{frame}{La ruta de aprendizaje}
  \centering
  \begin{tikzpicture}[node distance=0.42cm and 0.32cm,
    card/.style={rounded corners=4pt,minimum width=2.55cm,minimum height=2.25cm,
      text width=2.15cm,align=left,inner sep=8pt},
    arrow/.style={-{Latex[length=2mm]},line width=1pt,draw=SoftText!55}]
    \node[card,fill=Teal!12] (a) {\textcolor{Teal}{\bfseries 01}\\[0.2em]\bfseries Fundamentos\\[-0.1em]{\scriptsize Agente, entorno y retorno}};
    \node[card,fill=Coral!11,right=of a] (b) {\textcolor{Coral}{\bfseries 02}\\[0.2em]\bfseries Predicción\\[-0.1em]{\scriptsize Valores y métodos TD}};
    \node[card,fill=Gold!20,right=of b] (c) {\textcolor{Gold!70!Ink}{\bfseries 03}\\[0.2em]\bfseries Control\\[-0.1em]{\scriptsize SARSA y Q-learning}};
    \node[card,fill=Leaf!15,right=of c] (d) {\textcolor{Leaf!75!Ink}{\bfseries 04}\\[0.2em]\bfseries Deep RL\\[-0.1em]{\scriptsize DQN y policy gradients}};
    \draw[arrow] (a) -- (b); \draw[arrow] (b) -- (c); \draw[arrow] (c) -- (d);
  \end{tikzpicture}

  \vspace{0.55cm}
  \keyidea{Cada módulo conecta una ecuación, una intuición visual y un experimento reproducible.}
\end{frame}

\begin{frame}{Objetivos de aprendizaje}
  \begin{columns}[T,onlytextwidth]
    \column{0.48\textwidth}
      \begin{block}{Comprender}
        \begin{itemize}
          \item Formular problemas como MDP.
          \item Interpretar retorno, valor y ventaja.
          \item Explicar exploración vs. explotación.
        \end{itemize}
      \end{block}
    \column{0.48\textwidth}
      \begin{exampleblock}{Construir}
        \begin{itemize}
          \item Implementar algoritmos tabulares.
          \item Diagnosticar curvas de entrenamiento.
          \item Diseñar recompensas con criterio.
        \end{itemize}
      \end{exampleblock}
  \end{columns}

  \vspace{0.45cm}
  \begin{center}
    \state{$s_t$}\qquad $\longrightarrow$\qquad
    \rlaction{$a_t\sim\pi(\cdot\mid s_t)$}\qquad $\longrightarrow$\qquad
    \reward{$r_{t+1}$}
  \end{center}
\end{frame}

% ============================================================
\section{Fundamentos}
\module{Fundamentos}{Cómo aprende un agente cuando nadie le entrega la respuesta correcta}
% ============================================================

\begin{frame}{El bucle agente--entorno}
  \begin{columns}[c,onlytextwidth]
    \column{0.57\textwidth}
      \centering
      \begin{tikzpicture}[scale=0.98,transform shape,
        box/.style={rounded corners=5pt,minimum width=3.0cm,minimum height=1.35cm,
          font=\bfseries,align=center},
        flow/.style={-{Latex[length=2.5mm]},line width=1.25pt}]
        \node[box,fill=Teal,text=white] (agent) {AGENTE\\[-0.15em]\scriptsize política $\pi$};
        \node[box,fill=Ink,text=white,right=2.0cm of agent] (env) {ENTORNO\\[-0.15em]\scriptsize dinámica $p$};
        \draw[flow,Coral] ([yshift=0.28cm]agent.east) -- node[above,text=Ink,font=\small] {acción $a_t$} ([yshift=0.28cm]env.west);
        \draw[flow,Teal,bend left=28] (env.south) to node[below,text=Ink,font=\small,align=center] {estado $s_{t+1}$\\ recompensa $r_{t+1}$} (agent.south);
      \end{tikzpicture}
    \column{0.39\textwidth}
      \begin{itemize}
        \item No hay etiquetas de la acción correcta.
        \item Las decisiones cambian los datos futuros.
        \item La recompensa puede llegar con retraso.
      \end{itemize}
  \end{columns}

  \vspace{0.5cm}
  \keyidea{Aprender por refuerzo es aprender a \textbf{actuar} para maximizar recompensa acumulada.}
\end{frame}

\begin{frame}{Del problema al MDP}
  \begin{columns}[T,onlytextwidth]
    \column{0.46\textwidth}
      Un proceso de decisión de Markov se describe por
      \[
        \mathcal{M}=\langle \mathcal{S},\mathcal{A},p,r,\gamma\rangle.
      \]
      \begin{itemize}
        \item $\mathcal{S}$: estados observables
        \item $\mathcal{A}$: acciones disponibles
        \item $p$: dinámica del entorno
        \item $r$: señal de recompensa
        \item $\gamma$: preferencia temporal
      \end{itemize}
    \column{0.50\textwidth}
      \begin{block}{Caso conductor: Gridworld}
        \state{posición del agente}\par\smallskip
        \rlaction{arriba, abajo, izq., der.}\par\smallskip
        \reward{$-1$ por paso, $+10$ en la meta}
      \end{block}
      \vspace{0.3cm}
      \begin{alertblock}{Pregunta de diseño}
        ¿El estado contiene suficiente información para predecir el siguiente paso?
      \end{alertblock}
  \end{columns}
\end{frame}

\begin{frame}{Retorno: convertir el futuro en un número}
  \begin{columns}[c,onlytextwidth]
    \column{0.48\textwidth}
      \begin{beamercolorbox}[rounded=true,sep=1em]{block body}
        \centering
        {\Large $G_t=\displaystyle\sum_{k=0}^{\infty}\gamma^k r_{t+k+1}$}
      \end{beamercolorbox}
      \vspace{0.35cm}
      \begin{itemize}
        \item $\gamma\approx 0$: visión inmediata.
        \item $\gamma\approx 1$: visión a largo plazo.
        \item No siempre es solo un hiperparámetro: expresa el horizonte del problema.
      \end{itemize}
    \column{0.48\textwidth}
      \centering
      \begin{tikzpicture}[x=1.18cm,y=1cm]
        \draw[-{Latex},line width=1pt,SoftText] (0,0) -- (4.5,0);
        \foreach \x/\lab/\val/\c in {0/t/1/Teal,1/t+1/.72/Coral,2/t+2/.52/Gold,3/t+3/.37/Leaf,4/\cdots/.26/SoftText}{
          \fill[\c] (\x,0) circle (0.11);
          \draw[\c!70,line width=3pt] (\x,0.18) -- (\x,{0.18+1.35*\val});
          \node[below,font=\scriptsize] at (\x,-0.05) {$\lab$};
        }
        \node[align=center,font=\scriptsize,text=SoftText] at (2,-0.7) {el peso de las recompensas\\decrece con la distancia temporal};
      \end{tikzpicture}
  \end{columns}
\end{frame}

\begin{frame}{Ecuaciones de Bellman}
  \begin{columns}[T,onlytextwidth]
    \column{0.48\textwidth}
      \begin{block}{Valor de un estado}
        \[
          V^\pi(s)=\mathbb{E}_\pi\!\left[r_{t+1}+\gamma V^\pi(s_{t+1})\mid s_t=s\right]
        \]
      \end{block}
      \small ¿Qué retorno espero al comenzar en $s$ y seguir $\pi$?
    \column{0.48\textwidth}
      \begin{alertblock}{Valor de una acción}
        \[
          Q^\pi(s,a)=\mathbb{E}_\pi\!\left[r_{t+1}+\gamma Q^\pi(s_{t+1},a_{t+1})\right]
        \]
      \end{alertblock}
      \small ¿Qué retorno espero si primero elijo $a$ y luego sigo $\pi$?
  \end{columns}

  \vspace{0.55cm}
  \keyidea{Bellman transforma una predicción de largo plazo en \textbf{recompensa inmediata + predicción futura}.}
\end{frame}

% ============================================================
\section{Algoritmos}
\module{Algoritmos de control}{Aprender una política a partir de experiencia, paso a paso}
% ============================================================

\begin{frame}{Tres familias, una misma pregunta}
  \small
  \begin{tabular}{@{}p{2.35cm}p{3.55cm}p{3.25cm}p{2.5cm}@{}}
    \toprule
    \textbf{Método} & \textbf{Objetivo} & \textbf{Actualización} & \textbf{Rasgo} \\
    \midrule
    Monte Carlo & estimar desde episodios & retorno completo $G_t$ & sin sesgo, alta varianza \\
    \addlinespace
    SARSA & control on-policy & $r+\gamma Q(s',a')$ & aprende lo que ejecuta \\
    \addlinespace
    Q-learning & control off-policy & $r+\gamma\max_{a'}Q(s',a')$ & aprende una política greedy \\
    \bottomrule
  \end{tabular}

  \vspace{0.65cm}
  \begin{center}
    \begin{beamercolorbox}[rounded=true,sep=0.75em,wd=0.82\linewidth]{block body}
      \centering $Q(s,a)\leftarrow Q(s,a)+\alpha\big[\underbrace{\text{objetivo}-Q(s,a)}_{\text{error TD}}\big]$
    \end{beamercolorbox}
  \end{center}
\end{frame}

\begin{frame}[fragile]{Q-learning: plantilla de implementación}
  \begin{columns}[T,onlytextwidth]
    \column{0.61\textwidth}
      \begin{lstlisting}[style=rlcode]
for episode in range(n_episodes):
    state, _ = env.reset()
    done = False

    while not done:
        action = epsilon_greedy(Q[state], epsilon)
        next_state, reward, terminated, truncated, _ = env.step(action)
        done = terminated or truncated

        target = reward
        if not done:
            target += gamma * Q[next_state].max()

        Q[state, action] += alpha * (target - Q[state, action])
        state = next_state
      \end{lstlisting}
    \column{0.35\textwidth}
      \begin{block}{Leer el algoritmo}
        \begin{enumerate}
          \item Explorar.
          \item Observar transición.
          \item Construir objetivo TD.
          \item Corregir la estimación.
        \end{enumerate}
      \end{block}
      \vspace{0.2cm}
      \reward{el dato de aprendizaje}
  \end{columns}
\end{frame}

\begin{frame}{Exploración vs. explotación}
  \begin{columns}[c,onlytextwidth]
    \column{0.48\textwidth}
      \begin{block}{$\varepsilon$-greedy}
        \[
        a=\begin{cases}
          \text{acción aleatoria}, & \text{con prob. }\varepsilon,\\
          \arg\max_{a'}Q(s,a'), & \text{en otro caso.}
        \end{cases}
        \]
      \end{block}
      \small Una estrategia simple debe acompañarse de un calendario y una métrica de cobertura.
    \column{0.48\textwidth}
      \centering
      \begin{tikzpicture}[x=0.85cm,y=3.5cm]
        \draw[-{Latex},SoftText] (0,0) -- (5.5,0) node[right,font=\scriptsize] {episodio};
        \draw[-{Latex},SoftText] (0,0) -- (0,1.15) node[above,font=\scriptsize] {$\varepsilon$};
        \draw[Teal,line width=2.5pt,domain=0:5,samples=50] plot (\x,{0.1+0.9*exp(-0.65*\x)});
        \draw[Coral,dashed,line width=1.4pt] (0,0.1) -- (5,0.1);
        \node[font=\scriptsize,text=Teal] at (3.5,0.48) {decaimiento};
        \node[font=\scriptsize,text=Coral,anchor=west] at (3.8,0.15) {mínimo};
      \end{tikzpicture}
  \end{columns}

  \vspace{0.4cm}
  \keyidea{Reducir $\varepsilon$ demasiado pronto puede convertir una primera casualidad en una política permanente.}
\end{frame}

% ============================================================
\section{Laboratorio}
\module{Laboratorio y diagnóstico}{Entrenar es solo la mitad: hay que saber interpretar lo que ocurre}
% ============================================================

\begin{frame}{Cómo leer una curva de entrenamiento}
  \begin{columns}[c,onlytextwidth]
    \column{0.58\textwidth}
      \centering
      \begin{tikzpicture}[x=0.085cm,y=0.052cm]
        \draw[-{Latex},SoftText] (0,0) -- (75,0) node[right,font=\scriptsize] {episodios};
        \draw[-{Latex},SoftText] (0,0) -- (0,72) node[above,font=\scriptsize,align=center] {retorno\\medio};
        \fill[Teal!12] (0,10) .. controls (18,5) and (25,34) .. (38,29)
          .. controls (50,25) and (58,58) .. (72,51)
          -- (72,64) .. controls (58,70) and (48,40) .. (38,43)
          .. controls (24,48) and (18,18) .. (0,22) -- cycle;
        \draw[Teal,line width=2.2pt] (0,16) .. controls (18,11) and (25,41) .. (38,36)
          .. controls (50,32) and (58,65) .. (72,58);
        \draw[Coral,dashed,line width=1.2pt] (0,55) -- (72,55);
        \node[font=\scriptsize,text=Coral,anchor=west] at (48,51) {umbral objetivo};
      \end{tikzpicture}
    \column{0.38\textwidth}
      \begin{itemize}
        \item Mostrar media \emph{y} dispersión.
        \item Separar retorno de entrenamiento y evaluación.
        \item Repetir con varias semillas.
        \item Vigilar estabilidad, no solo máximo.
      \end{itemize}
  \end{columns}
\end{frame}

\begin{frame}{Ficha de experimento reproducible}
  \begin{columns}[T,onlytextwidth]
    \column{0.48\textwidth}
      \begin{block}{Configuración}
        \begin{itemize}
          \item Entorno y versión
          \item Semillas y presupuesto
          \item $\alpha$, $\gamma$, $\varepsilon$
          \item Criterio de terminación
        \end{itemize}
      \end{block}
    \column{0.48\textwidth}
      \begin{exampleblock}{Evidencia}
        \begin{itemize}
          \item Curvas con intervalos
          \item Política final visualizada
          \item Ablación de hiperparámetros
          \item Fallos y limitaciones
        \end{itemize}
      \end{exampleblock}
  \end{columns}

  \vspace{0.45cm}
  \keyidea{Una ejecución afortunada es una anécdota; varias semillas y un protocolo claro son evidencia.}
\end{frame}

\begin{frame}{Actividad: diseña el agente}
  \begin{columns}[T,onlytextwidth]
    \column{0.61\textwidth}
      \textbf{Escenario}\par
      Un sistema debe gestionar la climatización de un edificio reduciendo consumo sin sacrificar confort.

      \vspace{0.35cm}
      \begin{enumerate}
        \item Define $\mathcal{S}$, $\mathcal{A}$ y $r$.
        \item Decide si el problema es episódico o continuo.
        \item Identifica una posible recompensa mal especificada.
        \item Propón una política base para comparar.
      \end{enumerate}
    \column{0.35\textwidth}
      \begin{alertblock}{Criterio de éxito}
        La solución debe ser aprendible, medible y segura, no solo producir una recompensa alta.
      \end{alertblock}
      \vspace{0.3cm}
      \centering\Large\textcolor{Teal}{12 min}
  \end{columns}
\end{frame}

% ============================================================
\section{Cierre}
% ============================================================

\begin{frame}{Resumen del módulo}
  \begin{columns}[T,onlytextwidth]
    \column{0.31\textwidth}
      \begin{block}{1 · Formular}
        Estado, acción, dinámica, recompensa y horizonte.
      \end{block}
    \column{0.31\textwidth}
      \begin{alertblock}{2 · Aprender}
        Estimar valores a partir del error temporal.
      \end{alertblock}
    \column{0.31\textwidth}
      \begin{exampleblock}{3 · Evaluar}
        Comparar políticas con evidencia reproducible.
      \end{exampleblock}
  \end{columns}

  \vspace{0.65cm}
  \begin{center}
    {\Large\bfseries Interactuar $\rightarrow$ estimar $\rightarrow$ mejorar $\rightarrow$ repetir}
  \end{center}
\end{frame}

\begin{frame}[plain]
  \begin{tikzpicture}[remember picture,overlay]
    \fill[Paper] (current page.south west) rectangle (current page.north east);
    \fill[Teal] (current page.south west) rectangle ([yshift=0.13\paperheight]current page.south east);
    \draw[Coral,line width=3pt] ([xshift=-2.2cm,yshift=1.7cm]current page.center)
      -- ([xshift=2.2cm,yshift=1.7cm]current page.center);
    \node[text=Ink,font=\fontsize{27}{30}\selectfont\bfseries] at ([yshift=0.65cm]current page.center) {¿Preguntas?};
    \node[text=SoftText,font=\large] at ([yshift=-0.45cm]current page.center) {La mejor política comienza con una buena pregunta.};
    \node[text=white,font=\small] at ([yshift=0.08\paperheight]current page.south) {contacto@institucion.edu};
  \end{tikzpicture}
\end{frame}

\end{document}
